\section{Strong natural response and a corotational material driver}
\label{mr:section}

\paragraph{Scope and status.}
This section supplies author-checked component proofs, with independent
mathematical audit pending. It strengthens the fixed-regularization
calculus of Section~\ref{st:section}; it does not supply the arbitrary-data
critical time bound. In particular, a material derivative is not estimated
by simply deleting the transport term from \eqref{st:difference}.
Transport changes the metric in the gradient minimization problem. We
compute that change, including its contribution on the zero set, and keep
the nonlocal commutator in the shell-residual equation.
The one-form pullback mechanism is classical; see, for example,
Constantin--Iyer \cite[Lemmas~3.2--3.3]{ConstantinIyer2008}. Their stochastic
representation theorem is not an input here. Every deterministic
pullback and response assertion used below is derived explicitly, and no
novelty claim is made.

\subsection{The fixed weighted space and moving-metric response}

Fix $f\in L^3(\mathbb R^3;\mathbb R^3)$ and put $w=w(f)$,
$\rho=|w|$, $n=w/\rho$ where $\rho>0$, and $n=0$ elsewhere. Set
\[
 A(z)=|z|z,\qquad J(z)=|z|^{1/2}z,\qquad
\]
\[
 M=DA(w)=\rho(I+n\otimes n),\qquad
 N=DJ(w)=\rho^{1/2}(I+\tfrac12 n\otimes n).
\]
Both matrices are zero on $\{w=0\}$. Let $\mathcal H_w$ be
$L^2(\rho\,dx;\mathbb R^3)$ with inner product
$\langle a,b\rangle_w=\int a^{\mathsf T}Mb$; values on the zero
set are identified. Let $E_w$ be the closure of the image of
$\mathcal G_3$ in this Hilbert space, and let $P_w$ and $L_w=I-P_w$
be its orthogonal projections onto $E_w$ and $E_w^\perp$.
The notation $P_w$ does not denote the ordinary Leray projection.
We use
\begin{equation}\label{mr:matrix}
 \int\rho|a|^2\le\|a\|_w^2\le2\int\rho|a|^2,
 \qquad
 \|a\|_w^2\le\|Na\|_2^2\le\tfrac98\|a\|_w^2.
\end{equation}
Indeed, the last squared ratios are $1$ in transverse directions and
$9/8$ in the radial direction. H\"older gives
$\|a\|_w\le(2\|w\|_3)^{1/2}\|a\|_3$ for $a\in L^3$.
No distributional gradient representative is asserted for an arbitrary
element of $E_w$.

\begin{theorem}[Moving-metric response]\label{mr:response}
Let $h_\varepsilon\to h$ in $L^3$, and let measurable matrix fields satisfy
\[
 B_\varepsilon=I+\varepsilon G+o(\varepsilon)
       \quad\hbox{in }L^\infty,
 \qquad G\in L^\infty(\mathbb R^3;\mathbb R^{3\times3}).
\]
For sufficiently small nonzero $\varepsilon$, let $a_\varepsilon$
minimize
\[
 \frac13\int|B_\varepsilon^{-\mathsf T}a|^3
       \quad\hbox{over }f+\varepsilon h_\varepsilon+\mathcal G_3.
\]
Define, with zero values when $w=0$,
\[
 c=(I-\tfrac12n\otimes n)Gw,\qquad b=G^{\mathsf T}w+c,
 \qquad z=L_wh+P_wb.
\]
Then, for both signs of $\varepsilon$,
\begin{align}
 z_\varepsilon:=\frac{a_\varepsilon-w}{\varepsilon}
       &\longrightarrow z\quad\hbox{strongly in }\mathcal H_w,
       \label{mr:weighted-limit}\\
 \frac{\|a_\varepsilon-w\|_3^3}{\varepsilon^2}
       &\longrightarrow0,\label{mr:cubic-small}\\
 \frac{J(a_\varepsilon)-J(w)}{\varepsilon}
       &\longrightarrow Nz\quad\hbox{strongly in }L^2,\label{mr:natural-limit}\\
 \frac{J(B_\varepsilon^{-\mathsf T}a_\varepsilon)-J(w)}{\varepsilon}
       &\longrightarrow N(z-G^{\mathsf T}w)
          \quad\hbox{strongly in }L^2.\label{mr:physical-limit}
\end{align}
In particular, taking $B_\varepsilon=I$ proves the Hadamard derivative
\begin{equation}\label{mr:hadamard}
 D[J\circ w](f)[h]=NL_wh:L^3\longrightarrow L^2.
\end{equation}
This statement does not assert an unweighted $L^3$ derivative of $w$.
\end{theorem}

\begin{proof}
Uniform invertibility and boundedness of $B_\varepsilon$ follow from its
expansion. The functional is strictly convex, coercive and weakly lower
semicontinuous on the closed affine subspace of the reflexive space $L^3$.
It has a unique minimizer. Its stress
\[
 A_\varepsilon(v)=B_\varepsilon^{-1}
                   A(B_\varepsilon^{-\mathsf T}v)
\]
annihilates $\mathcal G_3$ at $a_\varepsilon$. The base stress $A(w)$
also annihilates that space. Uniformly for small $\varepsilon$,
\begin{align}
 (A_\varepsilon(a)-A_\varepsilon(v))\cdot(a-v)
   &\ge c_0(|a|+|v|)|a-v|^2,\label{mr:monotonicity}\\
 |A_\varepsilon(a)-A_\varepsilon(v)|
   &\le C_0(|a|+|v|)|a-v|.\label{mr:stress-lip}
\end{align}
Here $c_0>0$ and $C_0<\infty$ are independent of $a,v,x,\varepsilon$.
For the untransformed stress the first inequality follows from
\[
 (A(a)-A(v))\cdot(a-v)
 -\tfrac12(|a|+|v|)|a-v|^2
   =\tfrac12(|a|-|v|)^2(|a|+|v|)\ge0.
\]
Apply this with $B_\varepsilon^{-\mathsf T}a$ and
$B_\varepsilon^{-\mathsf T}v$ to obtain \eqref{mr:monotonicity}.
The second follows by integrating $DA$, whose norm is at most $2|v|$.

Write $d_\varepsilon=a_\varepsilon-w$. Comparing the stresses against
$d_\varepsilon-\varepsilon h_\varepsilon\in\mathcal G_3$ gives
\begin{align*}
 &\langle A_\varepsilon(a_\varepsilon)-A_\varepsilon(w),d_\varepsilon\rangle\\
 &=\varepsilon\langle A_\varepsilon(a_\varepsilon)
            -A_\varepsilon(w),h_\varepsilon\rangle
   -\langle A_\varepsilon(w)-A(w),
                         d_\varepsilon-\varepsilon h_\varepsilon\rangle.
\end{align*}
All these pairings are $L^{3/2}$--$L^3$. Put
$X=\int\rho|d_\varepsilon|^2$, $Y=\|d_\varepsilon\|_3^3$.
Using \eqref{mr:monotonicity}, \eqref{mr:stress-lip}, and
$|A_\varepsilon(w)-A(w)|\le C|\varepsilon|\rho^2$, the identity implies
\[
 X+Y\le C|\varepsilon|(X^{1/2}+Y^{2/3})+C\varepsilon^2.
\]
For example, the mixed terms are bounded using
$\int\rho|d_\varepsilon||h_\varepsilon|
\le\|w\|_3^{1/2}\|h_\varepsilon\|_3X^{1/2}$ and
$\int\rho^2|d_\varepsilon|\le\|w\|_3^{3/2}X^{1/2}$.
Constants may depend on bounded norms of $h_\varepsilon$, on $w$, and on
the matrix expansion, but not on $\varepsilon$. Young's inequality gives
\begin{equation}\label{mr:rough}
 \|d_\varepsilon\|_w^2+\|d_\varepsilon\|_3^3=O(\varepsilon^2),
 \qquad \|z_\varepsilon\|_3=O(|\varepsilon|^{-1/3}).
\end{equation}
The potentially divergent last norm must be retained in the next step.

The formula for $DA$ gives a global bound
$\|DA(a)-DA(v)\|_{\rm op}\le4|a-v|$: expand the difference of
$a\otimes a/|a|$ and $v\otimes v/|v|$, choosing $|a|\ge|v|$, and
bound its three terms by $|a-v|$, $|v||a-v|/|a|$ and the same quantity.
The scalar-matrix part contributes at most $|a-v|$.
Thus the stress Taylor remainder is quadratic. Expanding the two matrix
factors yields, for every $v\in L^3$,
\begin{equation}\label{mr:taylor-stress}
 T_\varepsilon(v):=\frac{A_\varepsilon(w+\varepsilon v)-A(w)}{\varepsilon}
       =M(v-b)+r_\varepsilon+e_\varepsilon(v),
\end{equation}
where pointwise
\[
 |r_\varepsilon|\le o(1)\rho^2,
 \qquad |e_\varepsilon(v)|
       \le C|\varepsilon|(\rho|v|+|v|^2).
\]
The $o(1)$ is uniform in space. The identity $Mc=GA(w)$ verifies the
constant term, without dividing by $\rho$ on its zero set.
The pairing of $r_\varepsilon$ against an $\mathcal H_w$-bounded family
tends to zero by weighted Cauchy--Schwarz. By \eqref{mr:rough},
$e_\varepsilon(z_\varepsilon)\to0$ in $L^{3/2}$.

The family $z_\varepsilon$ is bounded in $\mathcal H_w$ and lies in
$h_\varepsilon+E_w$. Every weak subsequential limit $z_*$ therefore
lies in $h+E_w$. Test \eqref{mr:taylor-stress} against each fixed
$g\in\mathcal G_3$. Stationarity and the remainder bounds imply
$\langle z_*-b,g\rangle_w=0$. Density extends this to $E_w$.
The unique possible limit is $z=L_wh+P_wb$, so the whole family converges
weakly to $z$.

To prove strong convergence and the essential extra limit
\eqref{mr:cubic-small}, choose $z_m=h+g_m\to z$ in $\mathcal H_w$
with $g_m\in\mathcal G_3$, and put
$z_{m,\varepsilon}=z_m+h_\varepsilon-h$. Fix $m$ before taking
$\varepsilon\to0$. Then
$z_\varepsilon-z_{m,\varepsilon}\in\mathcal G_3$, whence
\begin{align*}
 &\langle T_\varepsilon(z_\varepsilon)
                 -T_\varepsilon(z_{m,\varepsilon}),
                     z_\varepsilon-z_{m,\varepsilon}\rangle\\
 &\hspace{1cm}=-\langle T_\varepsilon(z_{m,\varepsilon}),
                     z_\varepsilon-z_{m,\varepsilon}\rangle.
\end{align*}
Let $v_\varepsilon=z_\varepsilon-z_{m,\varepsilon}$.
The left side is at least
\[
 c_0\int(|a_\varepsilon|+|w+\varepsilon z_{m,\varepsilon}|)
                    |v_\varepsilon|^2.
\]
It controls $c_0|\varepsilon|\|v_\varepsilon\|_3^3$. It also
controls $c_0\int\rho|v_\varepsilon|^2$ up to an error tending to zero,
since that error is at most
\[
 C|\varepsilon|\|z_{m,\varepsilon}\|_3
            (\|z_\varepsilon\|_3+\|z_{m,\varepsilon}\|_3)^2
              =O_m(|\varepsilon|^{1/3}).
\]
On the right side of the identity the $r_\varepsilon$ error tends to zero
in the weighted dual pairing. The quadratic part of
$e_\varepsilon(z_{m,\varepsilon})$ contributes
$O_m(|\varepsilon|^{2/3})$; its $\rho|z_{m,\varepsilon}|$ part tends
to zero by weighted Cauchy--Schwarz. The right side therefore tends to
\[
 -\langle z_m-b,z-z_m\rangle_w=\|z_m-z\|_w^2,
\]
because $z-b\perp E_w$. Consequently
\[
 \limsup_{\varepsilon\to0}\|z_\varepsilon-z_{m,\varepsilon}\|_w^2
   \le C\|z_m-z\|_w^2,
 \quad
 \limsup_{\varepsilon\to0}|\varepsilon|
                \|z_\varepsilon-z_{m,\varepsilon}\|_3^3
   \le C\|z_m-z\|_w^2.
\]
Triangle inequalities, followed by $m\to\infty$, give strong weighted
convergence and $|\varepsilon|\|z_\varepsilon\|_3^3\to0$.
No uniform $L^3$ bound on $z_m$ was assumed.

Finally $DJ$ is globally $1/2$-H\"older. To see this, when
$|a-v|\ge\frac12\max(|a|,|v|)$ bound both derivatives by
$\frac32\max(|a|,|v|)^{1/2}$. In the complementary case the segment
stays away from zero, and differentiating the displayed formula for $DJ$
bounds its variation by $C\max(|a|,|v|)^{-1/2}|a-v|$.
Both cases give $C|a-v|^{1/2}$. Integration along a segment yields
\[
 |J(w+d)-J(w)-Nd|\le C|d|^{3/2}.
\]
Its squared $L^2$ error divided by $\varepsilon^2$ tends to zero by
\eqref{mr:cubic-small}. This proves \eqref{mr:natural-limit}.
For $\widetilde d_\varepsilon=B_\varepsilon^{-\mathsf T}a_\varepsilon-w$,
\[
 \widetilde d_\varepsilon/\varepsilon
       \longrightarrow z-G^{\mathsf T}w\quad\hbox{in }\mathcal H_w,
 \qquad
 \|\widetilde d_\varepsilon\|_3^3/\varepsilon^2\longrightarrow0.
\]
These follow by expanding the matrix and using the bounded weighted norms
in \eqref{mr:rough}. The same Taylor argument proves
\eqref{mr:physical-limit}. If $w=0$ almost everywhere the weighted space
is the zero space, and the cubic little-$o$ argument still proves the
unweighted natural limits. This also covers nontrivial zero sets.
\end{proof}

\subsection{Removing transport and rigid rotation without losing strain}

On the selected original Navier--Stokes branch use the convention
$G_{ij}=\partial_j u_i$, and write
\[
 S=\tfrac12(G+G^{\mathsf T}),\qquad
 \Omega=\tfrac12(G-G^{\mathsf T}),\qquad
 \mathcal D_t=\partial_t+u\cdot\nabla.
\]
All assertions in the next theorem are on compact classical intervals;
none of its right sides is claimed to be uniform as $t\uparrow T_*$.

\begin{theorem}[Exact corotational response and its action]\label{mr:material}
Let $V=J(w(u))$ and define $U=\mathcal D_tV-\Omega V$. Then, as
strong $L^2$ derivatives at interior classical times,
\begin{equation}\label{mr:exact-driver}
 U=N\left[L_w(\nu\Delta u-Sw)+P_wc_S\right],
 \qquad c_S=(I-\tfrac12n\otimes n)Sw.
\end{equation}
In particular, with
\begin{equation}\label{mr:action}
 \mathfrak A(t)=\|L_w(\nu\Delta u-Sw)\|_w^2+\|P_wc_S\|_w^2,
\end{equation}
we have
\begin{align}
 \mathfrak A&=\|U\|_2^2-\tfrac19\|n\cdot U\|_2^2,
 &\mathfrak A\le\|U\|_2^2&\le\tfrac98\mathfrak A,\label{mr:two-sided}\\
 \|U\|_2^2&\le\tfrac92\nu^2\int\rho|\Delta u|^2
              +\tfrac{81}{16}\int\rho^3\|S\|_{\rm op}^2.
 &&\label{mr:action-bound}
\end{align}
The action is measurable and locally bounded in time. The equation has
neither a pressure-gradient driver nor a separate vorticity-amplitude
driver; it retains the full strain and diffusion response.
\end{theorem}

\begin{proof}
Fix a time $t$ and let $\Phi_s$ be the flow from $t$ to $t+s$ generated
by the actual velocity, with $B_s=D\Phi_s$. The compact-interval Sobolev
package makes $u$ and its spatial derivatives bounded and time continuous.
Thus $\Phi_s$ is a proper smooth diffeomorphism, $\det B_s=1$, and
$B_s=I+sG(t)+o(s)$ uniformly in space. Pullback preserves $\mathcal G_3$:
for a compactly supported smooth potential,
$B_s^{\mathsf T}(\nabla\psi)\circ\Phi_s=\nabla(\psi\circ\Phi_s)$;
boundedness on $L^3$ and density prove the assertion on its closure.
The inverse flow proves surjectivity. This uses unweighted $L^3$ density,
not an unproved weighted density statement.

The pulled-back representative
$a_s=B_s^{\mathsf T}w(u(t+s))\circ\Phi_s$ is consequently the
minimizer in Theorem~\ref{mr:response} for input
$f_s=B_s^{\mathsf T}u(t+s)\circ\Phi_s$. Indeed, change variables using
$\det B_s=1$ in its cubic integral. Differentiating the input gives
\[
 \frac d{ds}f_s
   =B_s^{\mathsf T}(\nu\Delta u)(t+s)\circ\Phi_s
       +\nabla\left[(|u|^2/2-p)(t+s)\circ\Phi_s\right].
\]
Here $G^{\mathsf T}u=\nabla(|u|^2/2)$, and the Navier--Stokes
material equation supplied $\nu\Delta u-\nabla p$.
Integrating in $s$ proves $f_s=f_0+s h_s+g_s$, with
$h_s\to\nu\Delta u(t)$ in $L^3$ and $g_s\in\mathcal G_3$.
The last term changes no affine coset. All the gradient terms are in
$L^3$ by the compact-interval regularity package.

Apply \eqref{mr:physical-limit} to
$V(t+s)\circ\Phi_s=J(B_s^{-\mathsf T}a_s)$. It yields
\[
 \mathcal D_tV
   =N\left[L_w(\nu\Delta u-G^{\mathsf T}w)
                  +P_w(I-\tfrac12n\otimes n)Gw\right].
\]
For completeness, the strong derivative of $V(t)$ itself exists by
\eqref{mr:hadamard} and $u\in C^1_tL^3_x$.
Since $V(t)\in H^1$, differentiation of the fixed function under the
smooth volume-preserving flow gives $u(t)\cdot\nabla V(t)$ in $L^2$.
Splitting the time increment into these two parts identifies the displayed
material derivative with $V_t+u\cdot\nabla V$.
Now $n\cdot\Omega w=0$, so the two rotational contributions combine to
$(L_w+P_w)\Omega w=\Omega w$, and $N\Omega w=\Omega V$.
This proves \eqref{mr:exact-driver}.

The two terms in brackets in \eqref{mr:exact-driver} are orthogonal in
$\mathcal H_w$. Inequality \eqref{mr:matrix} proves the two bounds in
\eqref{mr:two-sided}; diagonalizing the radial and transverse matrices
proves its identity, since the inverse squared radial ratio is $8/9$.
On the zero set $U=0$ as an $L^2$ field. The identity also establishes
measurability without asserting continuity of the moving projections.

For the last estimate write $Sw=\alpha n+\tau$, where $\tau\perp n$.
Contraction of the projections and the quadratic triangle inequality give
\[
 \mathfrak A\le2\|\nu\Delta u\|_w^2
                    +2\|Sw\|_w^2+\|c_S\|_w^2.
\]
The last two terms sum to
$\int\rho(\frac92\alpha^2+3|\tau|^2)
\le\frac92\int\rho|Sw|^2$.
Also $\|\nu\Delta u\|_w^2\le2\nu^2\int\rho|\Delta u|^2$.
Multiplying by $9/8$ gives \eqref{mr:action-bound}.
Its terms are locally bounded because $w\in L^3$, $\Delta u\in L^3$
and $S\in L^\infty$ with compact-interval bounds.
\end{proof}

\begin{remark}[Check against the accepted first-order balance]\label{mr:balance}
The response recovers rather than changes the known quotient balance.
Indeed, $Mw=2A(w)$ implies $w\perp E_w$ and $L_ww=w$. Using
$N V=\frac32 A(w)$ in \eqref{mr:exact-driver} gives
\[
 \frac23\langle V,U\rangle
   =\langle A(w),\nu\Delta u-Sw\rangle
   =-\nu D_Q-\int V^{\mathsf T}SV.
\]
Transport and skew rotation contribute zero to the derivative of
$\mathcal Q=\|V\|_2^2/3$. Thus the positive action cannot be substituted
for the signed strain work in the first-order equation.
\end{remark}

\subsection{A material residual law with the commutator retained}

Use the finite smooth features, fixed penalty $\eta>0$, coefficients $a$,
and residual $r$ of Section~\ref{st:section}. Explicitly,
$\phi_j=g_j(\rho)=h_j(V)$, where
$h_j(v)=g_j(|v|^{2/3})$ is smooth and globally Lipschitz, and
\[
 \Psi=\Pi(V\otimes V),\quad \chi=\mathcal T\Psi,\quad
 r=\chi-\sum_j a_j\phi_j,\quad
 R=\|r\|_2^2+\eta|a|^2.
\]
Here $\Pi$ is the symmetric trace-free projection and
$\mathcal T F=\sum_{i,j}R_iR_jF_{ij}$.
Define the scalar commutator field
\begin{equation}\label{mr:commutator}
 \mathcal C_u[V]=[u\cdot\nabla,\mathcal T]\Psi
   +\mathcal T\Pi(\Omega V\otimes V+V\otimes\Omega V),
\end{equation}
with $[u\cdot\nabla,\mathcal T]\Psi
=u\cdot\nabla(\mathcal T\Psi)-\mathcal T(u\cdot\nabla\Psi)$.
This notation includes both terms; neither is discarded or given a sign.

\begin{proposition}[Corotational finite-residual equation]\label{mr:residual}
On every compact classical interval, $\mathcal C_u[V]\in L^{3/2}$ at
almost every time, with locally bounded norm, and
\begin{align}
 R'={}&2\left\langle r,
   \mathcal T\Pi(U\otimes V+V\otimes U)+\mathcal C_u[V]\right\rangle
       -2\left\langle r,\sum_j a_jDh_j(V)U\right\rangle.
       \label{mr:residual-equation}
\end{align}
The first pairing is $L^3$--$L^{3/2}$ and the second is $L^2$--$L^2$.
Put $M_6=\|V\|_6$, $W_c=\|U\|_2$,
$A_3=(\sum_j\|\phi_j\|_3^2)^{1/2}$ and
$\ell=(\sum_j\|Dh_j\|_\infty^2)^{1/2}$. With $C_p$ the
$L^p$ operator norm of $\mathcal T$ on matrix fields,
\begin{align}
 |R'|\le{}&2\left(C_3\sqrt{2/3}\,M_6^2+A_3\sqrt{R/\eta}\right)
        \left(2C_{3/2}M_6W_c+\|\mathcal C_u[V]\|_{3/2}\right)
       \nonumber\\
       &+2\ell\eta^{-1/2}W_cR.\label{mr:residual-bound}
\end{align}
The identities $K=\langle\sigma,r\rangle$ and $N_\rho^2\le R$ from
the speed-shell construction remain unchanged.
\end{proposition}

\begin{proof}
The radial feature satisfies $Dh_j(V)\Omega V=0$. Therefore
$\mathcal D_t\phi_j=Dh_j(V)U$ in $L^2$, by the Sobolev chain rule.
The tensor product rule, in $L^{3/2}$, gives
\[
 \mathcal D_t\chi
     =\mathcal T\Pi(U\otimes V+V\otimes U)+\mathcal C_u[V].
\]
All terms in \eqref{mr:commutator} are indeed in $L^{3/2}$:
$V\in H^1\subset L^6$, so $\nabla\Psi\in L^{3/2}$ and
$\Psi\in L^{3/2}$ by interpolation of $V$ between $L^2$ and $L^6$.
The velocity and $\Omega$ are bounded on the interval, and the Riesz
operator is bounded on $L^{3/2}$. No endpoint commutator estimate has
been invoked.

The exact ridge envelope identity of Section~\ref{st:section} is
$R'=2\langle r,\chi_t\rangle
      -2\langle r,\sum_j a_j\phi_{j,t}\rangle$.
Replace the time derivatives by material derivatives.
The additional term is $-2\int r u\cdot\nabla r=0$.
To justify this on the whole space, observe that $r\in L^2\cap L^3$,
$\nabla\chi\in L^{3/2}$ and $\nabla\phi_j\in L^2$.
Thus $r\nabla r\in L^1$ and $r^2\in W^{1,1}$, by the truncated
Sobolev chain rule and passage to the limit. Integration against spatial
cutoffs uses $\nabla\cdot u=0$; the boundary term is bounded by
$C\|u\|_\infty\|r\|_2^2$ times the inverse cutoff radius and tends to
zero. This proves \eqref{mr:residual-equation}, including its dual spaces.

Finally, completing the ridge square gives
$\|r\|_2\le\sqrt R$ and $|a|\le\sqrt{R/\eta}$, while
\[
 \|r\|_3\le C_3\sqrt{2/3}\,M_6^2+A_3\sqrt{R/\eta},\qquad
 \|\mathcal T\Pi(U\otimes V+V\otimes U)\|_{3/2}
       \le2C_{3/2}M_6W_c.
\]
H\"older and $\|\sum_j a_jDh_j(V)U\|_2\le|a|\ell W_c$
prove \eqref{mr:residual-bound}.
\end{proof}

\begin{remark}[Rigid-motion covariance is only a diagnostic]\label{mr:rigid}
For a compactly supported smooth test field $V$, the commutator
\eqref{mr:commutator} vanishes when $u=b+\Omega x$ with constant $b$
and constant skew matrix $\Omega$. One way to verify this is to push $V$
forward by the corresponding rigid flow. The Fourier multiplier
$-\xi_i\xi_j/|\xi|^2$ contracts a rotated tensor into the rotated scalar,
so $\mathcal D_t(\mathcal T\Psi)=0$ when
$\mathcal D_tV=\Omega V$. The definition \eqref{mr:commutator} then
gives zero. A constant or rigidly rotating velocity is not an admissible
finite-energy whole-space datum; this verifies covariance, not an extra
Navier--Stokes regularity case. For general $u$, the commutator is retained.
\end{remark}

\paragraph{The first unsupported estimate.}
The stronger response theorem removes a zero-set differentiation obstacle,
and \eqref{mr:exact-driver} removes kinematic terms without treating them
as small. It does not close the missing bound. The available upper estimate
\eqref{mr:action-bound} contains
\[
 \nu^2\int\rho|\Delta u|^2,
       \qquad \int\rho^3\|S(u)\|_{\rm op}^2,
\]
which are not bounded at the endpoint by the energy or moment identities
proved earlier. The complete residual equation additionally requires
control of $\mathcal C_u[V]$, and \eqref{mr:residual-bound} still has
feature and $\eta^{-1/2}$ losses. The exact action, not its separate
absolute upper bounds, is a potential place to exploit cancellation.
Neither an input-only bound on the combined clock $L_c$ nor on
$\int_0^{\min(H,T_*)}\|\sigma\|_2^4\,dt$ has been proved here.
In particular, differentiating a moving weighted projection, assuming
heat convexity of the vector-valued cubic quotient, or omitting the
commutator would insert an additional unproved step. These component
proofs do not promote any open node or independently certify an earlier
review-pending supplement.
